\chapter{符号表与词汇表}

数学符号不是密码，而是缩写——每个符号背后都藏着一个“说来话长”的直观想法。本附录把全书出现的主要符号和术语按领域分类整理。表格里“直观解释”是你第一眼该想到的图景，“例子”是最小可用的验证。遇到不认识的符号，先查这里，再回到正文看它干活的样子。使用时有两类姿势：正文里遇到生词，按领域翻到对应小节查；想预习某一章，可以先把相关表格扫一遍，混个脸熟。表格不是背诵材料，而是地图——走得多了，路线自然就记住了。

\section{数}

数的世界是一层一层扩建的：为了做减法添了负数，为了做除法添了分数，为了写出“边长为 $1$ 的正方形的对角线”添了无理数。下表的符号就是这些楼层的名牌。

\begin{center}
\begin{tabular}{llp{4.6cm}p{4.6cm}}
\toprule
符号/术语 & 读法 & 直观解释 & 例子 \\
\midrule
$\mathbb{N}$ & 自然数集 & 数数用的数：$0$ 或 $1$ 起一个一个数下去 & $0,1,2,3,\dots$ \\
$\mathbb{Z}$ & 整数集 & 自然数加上它们的“欠账”（负数） & $-3,0,7$ \\
$\mathbb{Q}$ & 有理数集 & 能写成两整数之比的数 & $\dfrac{2}{3}$，$0.25$ \\
$\mathbb{R}$ & 实数集 & 数轴上所有的点，没有缺口 & $\sqrt{2}$，$\pi$ \\
$a\mid b$ & $a$ 整除 $b$ & $b$ 能被 $a$ 正好分完，没有剩余 & $3\mid 12$，但 $3\nmid 14$ \\
$\gcd(a,b)$ & 最大公约数 & 能同时整除 $a$ 和 $b$ 的最大整数 & $\gcd(12,18)=6$ \\
$|x|$ & 绝对值 & 数轴上 $x$ 到原点的距离，不分方向 & $|-5|=5$ \\
$\approx$ & 约等于 & 近似相等，误差在可接受范围内 & $\pi\approx 3.14$ \\
\bottomrule
\end{tabular}
\end{center}

\section{代数}

代数的核心动作是“用字母代替数”，于是运算规则可以脱离具体数字单独研究。字母 $x$ 不是“某个神秘的数”，而是一个占位符：它占着位置，让规律现身。

\begin{center}
\begin{tabular}{lp{4.8cm}p{5.4cm}}
\toprule
符号/术语 & 直观解释 & 例子 \\
\midrule
$a^n$ & $n$ 个 $a$ 连乘，乘法的缩写 & $2^5=32$ \\
$\sqrt{a}$ & 平方后等于 $a$ 的那个非负数 & $\sqrt{9}=3$ \\
$\log_a b$ & “$a$ 的几次方等于 $b$”的答案 & $\log_2 8=3$ \\
$\sum$ & 求和号：把一串数按规则全加起来 & $\sum_{k=1}^{4}k=1+2+3+4=10$ \\
$f(x)$ & 函数：一台机器，送进 $x$ 吐出唯一结果 & $f(x)=2x$，则 $f(3)=6$ \\
多项式 & 未知数只有加、减、乘的式子 & $x^2-3x+2$ \\
因式分解 & 把多项式拆成几个式子的乘积 & $x^2-1=(x+1)(x-1)$ \\
\bottomrule
\end{tabular}
\end{center}

\section{几何}

几何符号大多脱胎于图形本身：看到 $\triangle$ 想到三角形，看到 $\perp$ 想到墙角。记符号时顺手在纸上画一下，比默背十遍更有效。

\begin{center}
\begin{tabular}{lp{4.8cm}p{5.4cm}}
\toprule
符号/术语 & 直观解释 & 例子 \\
\midrule
$\angle ABC$ & 顶点在 $B$ 处张开的角度 & $\angle ABC=60^\circ$ \\
$\triangle ABC$ & 以 $A,B,C$ 为顶点的三角形 & 三角形内角和 $180^\circ$ \\
$\cong$ & 全等：能完全重合 & 边边边判全等 \\
$\sim$ & 相似：形状相同，大小可不同 & 放大照片与原照片相似 \\
$\perp$ & 垂直：夹角恰好 $90^\circ$ & 墙角的两边互相垂直 \\
$\parallel$ & 平行：同一平面内永不相交 & 铁轨的两条钢轨 \\
$\pi$ & 圆周与直径的比，对所有圆都一样 & 周长 $=2\pi r$ \\
\bottomrule
\end{tabular}
\end{center}

\section{分析}

分析（微积分的严格化）研究“无限逼近”，它的符号几乎都在回答同一个问题：离目标到底有多近？把“近”说精确的艰难过程，正是正文第 $9$ 到 $12$ 章的故事。

\begin{center}
\begin{tabular}{lp{4.8cm}p{5.4cm}}
\toprule
符号/术语 & 直观解释 & 例子 \\
\midrule
$\lim$ & 极限：无限逼近但不必到达的目标 & $\dfrac{1}{n}\to 0$（$n$ 越大越接近 $0$） \\
$\infty$ & 无穷大：比任何数都大的“方向”，不是数 & $1,2,3,\dots$ 数下去没有尽头 \\
$\varepsilon$ & 任意小的误差预算，你挑多小都行 & 精确到 $\varepsilon=0.001$ \\
$\Delta x$ & 变化量：$x$ 改动的一小步 & 从 $3$ 到 $3.1$，$\Delta x=0.1$ \\
$f'(x)$ & 导数：曲线在一点的瞬时坡度 & 位移的导数是速度 \\
$\int$ & 积分：把无穷多个小条加起来求总量 & 速度对时间积分得路程 \\
\bottomrule
\end{tabular}
\end{center}

\section{概率}

概率论给“不确定”立规矩：可能性用 $0$ 到 $1$ 之间的数表示，$0$ 是不可能，$1$ 是必然。它的难点往往不在计算，而在把现实问题翻译成正确的“事件”。

\begin{center}
\begin{tabular}{lp{4.8cm}p{5.4cm}}
\toprule
符号/术语 & 直观解释 & 例子 \\
\midrule
$P(A)$ & 事件 $A$ 发生的可能性，$0$ 到 $1$ 之间 & 掷硬币 $P(\text{正面})=\dfrac{1}{2}$ \\
$P(A\mid B)$ & 已知 $B$ 发生后，$A$ 发生的概率 & 已知下雨，带伞的概率 \\
$E(X)$ & 期望：长期平均下来每次得到多少 & 掷骰子的期望是 $3.5$ \\
随机变量 & 结果带随机性的“测量值” & 掷骰子得到的点数 \\
方差 & 数据围绕期望的波动大小 & 成绩越稳定，方差越小 \\
\bottomrule
\end{tabular}
\end{center}

\section{离散数学}

离散数学研究“一个一个分开”的对象：整数、集合、图。它是计算机科学的母语——计算机里的东西没有“半个”，全是离散的。附录 A 里的量词符号也收录在本节，它们是全书写证明时最频繁打交道的几个记号。

\begin{center}
\begin{tabular}{lp{4.8cm}p{5.4cm}}
\toprule
符号/术语 & 直观解释 & 例子 \\
\midrule
$\in$ & 属于：某个东西在这个集合里 & $2\in\{1,2,3\}$ \\
$\forall$ & 任意：每一个都要满足 & 所有偶数能被 $2$ 整除 \\
$\exists$ & 存在：至少能找到 & 存在一个偶质数 \\
$a\equiv b\pmod{m}$ & 同余：$a$ 和 $b$ 除以 $m$ 余数相同 & $17\equiv 2\pmod{5}$ \\
$n!$ & 阶乘：$n$ 个东西排成一列的排法数 & $3!=6$ \\
$\dbinom{n}{k}$ & 组合数：从 $n$ 个中无序挑 $k$ 个的挑法数 & $\dbinom{4}{2}=6$ \\
图（graph） & 一些点和点之间的连线组成的网络 & 城市为点、公路为边的交通图 \\
\bottomrule
\end{tabular}
\end{center}

一个读符号的通用建议：遇到陌生符号，先大声读出来，再问自己三件事——它作用在什么对象上？它回答什么问题？它的“反操作”是什么？比如 $\log$ 作用在正数上，回答“几次方”，反操作是指数。三件事答齐了，符号就不再吓人。

\begin{lianjie}
符号是会“串门”的：$\sum$ 在代数里求有限和，到了分析里变成无穷级数；$\in$ 在集合里登记成员，到了概率里又用来写“事件包含某个结果”。读任何一章遇到符号发懵时，先回本附录对号入座，再看它在新语境里被赋予了什么新任务。
\end{lianjie}

最后提醒：一个符号的含义由它的定义决定，而不是由它的长相决定。同一个字母在不同章节可能扮演不同角色——$x$ 在方程里是未知数，在函数里是自变量，在几何里可能是横坐标。每次见面先问一句“这次它是谁”，是读数学书的基本功。
